\begin{solapartat}
\figura[0.4]{sol_fig1.png}

\punts{0,25} Esquema: han d'aparèixer les plaques metàl·liques horitzontals, el voltatge i les dues forces que actuen sobre l'electró.

\punts{0,25} Les dues plaques formen un condensador pla. La placa superior ha de tenir una càrrega positiva i la inferior negativa per crear un camp elèctric vertical i dirigit cap avall de manera que creï una força vertical cap amunt en la càrrega negativa de la gota per contrarestar el pes dirigit cap avall.

\punts{0,5} El potencial elèctric és l'integral del camp elèctric i en aquest cas on el camp és homogeni, $|\vec{E}| = \frac{V}{d} = \frac{2000}{0,02} = 10^{5}\ \mathrm{V/m}$.

\punts{0,25} El camp elèctric és homogeni i dirigit cap avall; per tant, les línies de camp seran línies verticals equiespaiades i cap avall.
\end{solapartat}

\begin{solapartat}
\punts{0,75} El pes de la gota és $P = m\cdot g$, i la massa la calculem amb la densitat i el volum de la gota esfèrica: $m = \rho\cdot\frac{4}{3}\pi r^3 = 923\cdot\frac{4}{3}\pi(1,08\cdot 10^{-6})^3 = 4,87\cdot 10^{-15}\ \mathrm{kg}$.

La força elèctrica és $|\vec{F}_e| = q|E|$.

En l'equilibri la força elèctrica sobre la gota ha de ser igual al seu pes: $F_e = P$.

Igualant les dues forces: $|\vec{F}_e| = P = m\cdot g$, per tant $q = \frac{m\cdot g}{E} = \frac{4,87\cdot 10^{-15}\cdot 9,8}{10^{5}} = 4,77\cdot 10^{-19}\ \mathrm{C}$

\punts{0,25} Dividint la càrrega per la càrrega de l'electró $|e|$ trobem el número d'electrons $n$:
\[ n = \frac{4,77\cdot 10^{-19}\ \mathrm{C}}{1,602\cdot 10^{-19}\ \mathrm{C}} = 3 \]

\punts{0,25} Si la gota perd un electró, el mòdul de la força elèctrica esdevé menor que el pes, i la gota cau: $|\vec{F}_e| = q|E| < P$.

Alternativament, l'alumne/a ho pot justificar calculant la força elèctrica per dos electrons i el pes: $|\vec{F}_e| = q|E| = 3,20\cdot 10^{-14}\mathrm{N}$, $P = mg = 4,87\cdot 10^{-15}\cdot 9,8 = 4,77\cdot 10^{-14}\ \mathrm{N}$. I, per tant, $|\vec{F}_e| < P$, i la gota cau.
\end{solapartat}
